\section{Introduction}
An agent Skill packages procedural instructions with scripts and supporting resources. It can help a language-model agent complete a task consistently, but it also places reusable third-party instructions close to an agent's decision process. A package can request a plausible action that the current user never authorized. A backup step may disclose a private record; an update script may overwrite unrelated metadata; an integration check may contact an additional service. These actions need not look intrinsically malicious. Their permissibility depends on the current task, the data involved, and the destination of the resulting effect.

Consider a user who requests that an agent copy the summary of an input record into an existing report and send that input record to an approved service. The user permits local inspection of a private diagnostic file but prohibits its disclosure. A Skill asks the agent to bundle the diagnostic file with the approved upload. A destination allowlist alone cannot distinguish the legitimate upload from the larger disclosure: both contact the same service. A whole-Skill refusal prevents disclosure but may also prevent the requested report. A more selective response removes the unauthorized data while preserving the permitted update and upload. If the data have already been combined into an opaque value, however, the selective response may be unable to separate them without losing useful content.

This example exposes three different questions. First, what authority follows from the user request? Second, what actions and data paths does the Skill propose? Third, which restrictions prevent a forbidden effect while preserving the requested result? Treating all three as a single language-model judgment makes it difficult to identify whether a failure came from the authority interpretation, the reconstruction of the Skill, or the enforcement mechanism. Separating them creates inspectable intermediate artifacts and enables stronger experimental comparisons.

Existing research already addresses substantial parts of this problem. Skill-Inject evaluates contextual and explicit attacks placed in Skill packages~\cite{skillinject2026}. SkillScope analyzes actions relative to tasks and constrains validated over-privilege~\cite{skillscope2026}. SkillGuard connects permission declarations to context and action governance~\cite{skillguard2026}. Progent generates and deterministically mediates task-dependent tool policies~\cite{progent2025}; CaMeL tracks capabilities and information flow~\cite{camel2025}. These studies make it inappropriate to claim that task-aware permission control, graph-based Skill analysis, or runtime mediation is new by itself. The narrower question investigated here is whether a goal-path analysis adds practical value to selective repair when a strong task-only gateway is already available.

We implement \emph{goal-aware capability reconciliation} (\method), which independently extracts a task contract and a Skill action graph. A deterministic procedure reconciles potential paths with the contract and selects candidate restrictions at implicated sinks. A closed runtime then enforces the contract and, where selected, filters unauthorized data handles or write fields. The implementation operates over structured records and simulated external deliveries. It neither executes arbitrary upstream scripts nor implements operating-system confinement for an unrestricted coding agent. This scope makes every relevant effect observable, but limits the conclusions that can be drawn about deployed agents.

Our research questions are:
\begin{enumerate}
\item \textbf{RQ1: Representation.} How accurately do task contracts and Skill graphs recover the finite permissions and observable data dependencies in controlled workflows, and how do static judgments agree with original benchmark labels?
\item \textbf{RQ2: Security and utility.} How do coarse blocking, task-only mediation, and goal-targeted repair compare in observed unauthorized effects, task completion, false blocking, and joint success?
\item \textbf{RQ3: Incremental mechanism.} Does goal-targeted repair improve on an equally selective task-only repair gateway, and what changes when task authority, provenance, or enforcement is removed?
\end{enumerate}

Before the held-out execution, we specified directional hypotheses: source-aware enforcement should outperform destination-only checks under forced mixed-data transfers; selective repair should preserve more useful work than whole-Skill blocking; and goal analysis should be judged against an equally capable task-only repair mechanism. The third hypothesis is deliberately demanding. A benefit over a coarse or less expressive baseline would not establish that graph reasoning caused the benefit. Conversely, equal outcomes would support the simpler explanation that task interpretation and the runtime gate dominate in this setting.

The study uses \luna{} for every experimental model role and distinguishes three evidence sources: original Skill-Inject bundles for static analysis, authored structured workflows for controlled execution, and original attack text transferred into new fixtures. It further separates live model proposals from deterministic replay of the supplied workflow. Replay exposes the gateway to a known action sequence; it does not measure whether a language model would choose that sequence. We preserve development records, freeze formal inputs, and evaluate outcomes with deterministic checks independent of the defense matcher. No independent human annotation is claimed.

The contribution is consequently an auditable implementation and a controlled comparison, rather than a general security theorem or a native benchmark reproduction. The released materials connect each outcome to its task, bundle, model response, contract, graph, runtime events, and acceptance criteria. This structure also permits negative findings: if an undefended model avoids the tested injections, or a task-only repair method matches the proposed pipeline, those observations constrain the claim instead of being removed from the experiment.
